%!TEX root = hems_proj.tex
%%%%%%%%%%%%%%%%%%%%%%%%%%%%%%%%%%%%%%%%%%%%%%%%%%%%%%%%%%%%%%%%%%%%%%%
\chapter{Reprodukálhatóság és Érvényességi Korlátok}\label{ch:REPRO}
%%%%%%%%%%%%%%%%%%%%%%%%%%%%%%%%%%%%%%%%%%%%%%%%%%%%%%%%%%%%%%%%%%%%%%%

\begin{osszefoglal}
A fejezet célja a projekt tudományos megbízhatóságának dokumentálása: hogyan reprodukálhatók a futások, mely tényezők torzíthatják az eredményeket, és milyen bővítések szükségesek a külső validitás növeléséhez.
\end{osszefoglal}

\section{Reprodukálhatósági alapelvek}
A jelenlegi benchmark reprodukálhatóságát három technikai elv biztosítja:
\begin{enumerate}
  \item determinisztikus seedelési séma,
  \item fixen deklarált naplista és algoritmuskészlet,
  \item strukturált output mentés nyers és aggregált szinteken.
\end{enumerate}

A seedelés profile-, algorithm- és run-szinten eltolt, így egyszerre támogatja a megismételhetőséget és a kontrollált variabilitást.

\section{Futási környezet és artefaktok}
A projekt kimeneteit célszerű három szinten kezelni:
\begin{itemize}
  \item \textbf{Nyers szint}: minden futás külön sorban (költség, futásidő);
  \item \textbf{Aggregált szint}: algoritmusonkénti átlag/szórás/gap/győzelem;
  \item \textbf{Narratív szint}: markdown report és dolgozati ábrák.
\end{itemize}

Ezzel biztosítható, hogy a dolgozat szöveges állításai közvetlenül ellenőrizhetők legyenek táblázatos formában.

\section{Belső érvényesség (internal validity)}
A belső érvényességet főként az alábbi tényezők befolyásolják:
\begin{itemize}
  \item algoritmusparaméterek (populációméret, iterációszám),
  \item futások száma profilonként,
  \item seed-függő sztochasztikus szórás,
  \item korlátbüntetések súlyozása a célfüggvényben.
\end{itemize}

A jelen konfiguráció gyors összehasonlítást céloz, ezért az eredmények rangsor-jellegű értelmezése megbízhatóbb, mint a finom abszolút különbségek túlinterpretálása.

\section{Külső érvényesség (external validity)}
A külső validitás jelenleg közepes erősségű, mivel:
\begin{itemize}
  \item a PV adatok valósak,
  \item az árprofilok valós ENTSO-E forrásból származnak,
  \item ugyanakkor a load profil döntően generált.
\end{itemize}

A gyakorlati általánosíthatóság javításához szükséges:
\begin{enumerate}
  \item mért háztartási load-idősorok bevonása,
  \item több háztartás- és eszközkonfiguráció vizsgálata,
  \item hosszabb időhorizontú (szezonális/éves) validáció.
\end{enumerate}

\section{Fenyegetések az érvényességre (threats to validity)}
\subsection{Adatminőségi kockázatok}
Időbélyeg-inkonzisztencia, hiányos órák, forrásoldali formátumváltás mind torzíthatják a napi profilképzést. A fallback logika csökkenti a futásmegszakadás esélyét, de elemzési torzítást okozhat, ha túl gyakran aktiválódik.

\subsection{Modellezési egyszerűsítések}
A jelen célfüggvény még nem tartalmaz explicit akkumulátor-degradációs költséget, és a visszatáplálási ár egyszerűsített szorzóval modellezett. Ezek érdemi közelítések, de korlátozzák a gazdasági becslések pontosságát.

\subsection{Algoritmikus torzítások}
Kis futásszám mellett egy-egy szerencsés/unlucky seed túlreprezentálódhat. Ezért a jövőbeli, végleges összehasonlításban a futásszám növelése és bootstrap-alapú konfidencia-intervallumok alkalmazása ajánlott.

\section{Javasolt validációs protokoll a következő iterációra}
A következő fejlesztési körben célszerű standardizált validációs protokollt használni:
\begin{enumerate}
  \item baseline definíció (akku és ütemezés nélkül),
  \item fix train/validation naplista,
  \item legalább 10--20 futás algoritmus/profil szinten,
  \item összesített statisztika konfidencia-intervallummal,
  \item érzékenységvizsgálat (ár, degradáció, eszközparaméterek).
\end{enumerate}

E protokoll alkalmazásával a dolgozat eredményei nemcsak reprodukálhatók, hanem erősebb tudományos alátámasztást is kapnak.
